\subsection*{NVMe Components}

NVMe (Non-Volatile Memory Express) consists of hardware, software, and system-level components that work together to provide high-speed communication between the operating system and solid-state drives.

\subsubsection*{Hardware Components}

\begin{table}[H]
\begin{tabular}{|l|l|}
\hline
Component &
  Description \\ \hline
NVMe Controller &
  \begin{tabular}[c]{@{}l@{}}The controller is the main processing unit of the SSD. It executes\\ commands, manages data transfers, performs error correction,\\ and coordinates communication between the host system and\\ NAND flash memory.\end{tabular} \\ \hline
NAND Flash Memory &
  \begin{tabular}[c]{@{}l@{}}The primary non-volatile storage medium used to permanently\\ store user data. Common types include SLC, MLC, TLC, and QLC,\\ each providing different levels of speed, endurance, and storage\\ capacity.\end{tabular} \\ \hline
DRAM Cache &
  \begin{tabular}[c]{@{}l@{}}A high-speed memory buffer that stores frequently accessed data\\ and logical-to-physical address mapping tables, improving SSD\\ performance and reducing access latency.\end{tabular} \\ \hline
PCIe Interface &
  \begin{tabular}[c]{@{}l@{}}Provides a high-bandwidth communication link between the SSD\\ and the motherboard, enabling significantly faster data transfer\\ than the SATA interface.\end{tabular} \\ \hline
Power Management Unit &
  \begin{tabular}[c]{@{}l@{}}Regulates power consumption, supports multiple power-saving\\ states, and protects stored data during unexpected power\\ interruptions.\end{tabular} \\ \hline
Form Factor &
  \begin{tabular}[c]{@{}l@{}}Represents the physical design of the SSD. Common form factors\\ include M.2, U.2, and Add-in Card (AIC), each designed for different\\ computer systems.\end{tabular} \\ \hline
\end{tabular}
\caption{Hardware Components of NVMe}
\label{tab:hardware table}
\end{table}

\subsubsection*{Software and Protocol Components}

\begin{table}[H]
\begin{tabular}{|l|l|}
\hline
Component &
  Description \\ \hline
Submission Queue &
  \begin{tabular}[c]{@{}l@{}}Stores I/O commands issued by the operating system before they are\\ processed by the NVMe controller. Multiple queues enable highly\\ parallel command execution.\end{tabular} \\ \hline
Completion Queue &
  \begin{tabular}[c]{@{}l@{}}Returns the execution status and results of completed commands\\ back to the operating system.\end{tabular} \\ \hline
NVMe Commands &
  \begin{tabular}[c]{@{}l@{}}A standardized command set consisting of administrative commands\\ for device management and I/O commands for reading, writing, and\\ flushing data.\end{tabular} \\ \hline
Registers &
  \begin{tabular}[c]{@{}l@{}}Memory-mapped control and status registers that allow the operating\\ system to configure, monitor, and communicate with the NVMe\\ controller.\end{tabular} \\ \hline
Namespaces &
  \begin{tabular}[c]{@{}l@{}}Logical storage partitions that allow a single NVMe drive to present one\\ or more independent storage devices to the operating system.\end{tabular} \\ \hline
Firmware &
  \begin{tabular}[c]{@{}l@{}}Embedded software that controls NAND management, wear leveling,\\ garbage collection, bad block management, and error correction.\end{tabular} \\ \hline
\end{tabular}
\caption{Software Components of NVMe}
\label{tab:software table}
\end{table}

\subsubsection*{System-Level Components}

\begin{table}[H]
\begin{tabular}{|l|l|}
\hline
Component & Description \\ \hline
NVMe Driver &
  \begin{tabular}[c]{@{}l@{}}Operating system software that initializes the NVMe device, creates\\ queues, manages interrupts, and translates storage requests into\\ NVMe commands.\end{tabular} \\ \hline
File System Interface &
  \begin{tabular}[c]{@{}l@{}}Provides a standard block device abstraction that enables\\ applications to access NVMe storage using common file systems such\\ as NTFS, ext4, and APFS.\end{tabular} \\ \hline
DMA Engine &
  \begin{tabular}[c]{@{}l@{}}Transfers data directly between the SSD and system memory without\\ continuous CPU involvement, reducing processor overhead and\\ improving overall system performance.\end{tabular} \\ \hline
\end{tabular}
\caption{System-Level Components of NVMe}
\label{tab:system-level table}
\end{table}